\section{Сокращение моделей и численная точность}

\subsection{Сокращение модели и изменение арифметики}

Одним из ранних способов уменьшить потребность нейронных сетей в
вычислительных ресурсах стало удаление параметров.
В работе \textit{Optimal Brain Damage},
представленной Я.~Лекуном, Дж.~Денкером и С.~Соллой на конференции
1989 года, исследовалось удаление параметров на основании оценки
изменения функции ошибки~\cite{obd1989}. Предлагалось учитывать
значимость связи для поведения сети, а не только её наличие
в исходной архитектуре.

Часть связей может быть избыточной для решения выбранной задачи.
Удаление параметров изменяет набор
связей, тогда как сокращение разрядности изменяет множество
представимых значений. Указанные преобразования требуют различных
средств аппаратной реализации. Обработка нерегулярно разреженных
матриц предполагает организацию доступа к ненулевым элементам
и учёт связанных с этим управляющих затрат.

В работе Ю.~Гуна и соавторов 2014 года сжатие параметров свёрточных
сетей исследовалось средствами векторного квантования~\cite{gong2014}.
Вместо независимого хранения каждого исходного значения использовалось
представление через ограниченный набор кодов. Методы компактного
представления информации применялись к параметрам обученных моделей.
Авторы прежде всего оценивали сокращение памяти для хранения параметров.
Связь объёма сжатого представления со стоимостью матричной арифметики
зависит от способа декодирования и организации вычислений.

С.~Хань, Х.~Мао и У.~Дэлли объединили прореживание, обучаемое
квантование и кодирование Хаффмана в методе \textit{Deep Compression},
опубликованном в электронной форме в 2015 году и представленном
на Международной конференции по обучению представлений
(International Conference on Learning Representations, ICLR)
в 2016 году~\cite{han2016}. Совместное применение этих
этапов показало, что избыточность существует на нескольких уровнях:
в числе связей, в наборе значений весов и в способе хранения кодов.

Совместное применение сжатия и прореживания выявило необходимость
согласования представления модели с её исполнением. Сжатые коды
могут требовать декодирования, а представление разреженных связей —
хранения индексов и дополнительных обращений к памяти.
При оценке затрат на исполнение сжатой модели необходимо учитывать эти
дополнительные действия. Уменьшение объёма хранения и сокращение вычислений
проверяются отдельно, поскольку одно не всегда приводит к другому.


\subsection{Компактные архитектуры и дистилляция}

Затраты инференса можно сокращать, выбирая менее сложную модель, которая
удовлетворяет требованиям к качеству решения задачи. Наряду со сжатием уже
обученных сетей развивались методы обучения моделей с меньшим числом операций.
Ограничения памяти и вычислительных ресурсов в этом случае учитываются при
выборе архитектуры и способа обучения сети.

В работе Дж.~Хинтона, О.~Виньялса и Дж.~Дина 2015 года
развивался метод дистилляции знаний~\cite{hinton2015distillation}.
Результаты более сложной модели или ансамбля используются
при обучении менее сложной модели для последующего инференса. В задаче
классификации распределение ответов исходной модели даёт дополнительную
информацию об отношениях между классами. Таким образом, при обучении можно
использовать более сложную систему, чем та, которая будет затем обрабатывать
входные данные.

В MobileNets, предложенных Э.~Ховардом и соавторами в 2017 году,
использовались разделимые свёртки и параметры изменения ширины
сети и разрешения входа~\cite{howard2017mobilenets}.
Пространственная обработка отдельных каналов разделялась
с их последующим объединением посредством свёртки размером $1\times1$.
Архитектура предназначалась для мобильных и встраиваемых
приложений с ограниченными ресурсами. Изменение её параметров
позволяло сопоставлять качество решения задачи с задержкой
и размером модели в исследованных условиях.

Дистилляция и проектирование компактной архитектуры решают разные задачи при
подготовке модели. Дистилляция изменяет способ обучения с помощью ответов
другой модели, а выбор компактной архитектуры — структуру вычисляемых
преобразований. Оба подхода допускают последующее
квантование или аппаратную специализацию. Вместе с тем
сокращение числа операций требует отдельной оценки энергии:
результат зависит также от обмена данными, загрузки устройства
и поддержки выбранных операций исполнительной средой.

\subsection{Квантование и критерий достаточной точности}

Квантование заменяет исходное множество числовых значений конечным
набором уровней. В нейросетях допустимость такой замены проверяют
по качеству решения прикладной задачи.
В работе В.~Суна, С.~Сина и К.~Хвана 2015--2016 годов показана
зависимость устойчивости к квантованию от сложности сети
и возможности повторного обучения~\cite{sung2016}.

Устойчивость к квантованию определяется свойствами модели и условиями
её подготовки. Сеть с ограниченной выразительной способностью
и избыточно параметризованная модель могут по-разному реагировать
на одинаковые ограничения представления параметров. Оценка
допустимой погрешности требует совместного рассмотрения архитектуры
сети, процедуры обучения и численного формата. Разрядность операндов
характеризует представление, но не задаёт качество решения целевой задачи.

В бинарных нейронных сетях М.~Курбарио и соавторов 2016 года
веса и активации при исполнении ограничивались двумя значениями.
Это позволяло заменить значительную часть обычной арифметики
битовыми операциями~\cite{bnn2016}. При таком представлении менялся и набор операций, выполняемых аппаратурой.
Обучение с бинарными ограничениями формирует параметры с учётом
этого представления. Полученная модель отличается по процедуре
подготовки от результата непосредственного однобитного квантования
сети, обученной без соответствующих ограничений.

Применимость бинаризации определяется требуемым качеством решения
задачи, процедурой подготовки модели и поддерживаемыми аппаратными
операциями. Сокращение разрядности сопровождалось разработкой
методов компенсации численных изменений. Для каждой модели требовалось установить, насколько можно ограничить
представление чисел при заданных условиях её применения.

\subsection{Целочисленный инференс и смешанная точность}

В работе Б.~Якоба и соавторов, опубликованной в 2017 году
и представленной на конференции по компьютерному зрению и распознаванию
образов (Conference on Computer Vision and Pattern Recognition, CVPR)
в 2018 году, квантование было согласовано
с целочисленной реализацией инференса~\cite{jacob2018}.
Числовое представление описывалось масштабом и значением,
соответствующим нулю. Подготовка сети учитывала будущую
квантованную арифметику, а качество проверялось на прикладных задачах.

Целочисленный инференс предполагает согласование умножения,
накопления и преобразования результата к представлению следующего
слоя. Компактное хранение весов составляет лишь один из элементов
такой реализации. Форматы входных операндов, частичных сумм и выходов выбираются с учётом их
роли в вычислении.
В частности, разрядность накопления может отличаться
от разрядности множителей.

Другим подходом стало использование смешанной точности. В работе
П.~Мицкявичюса и соавторов 2017--2018 годов рассматривалось
обучение с использованием 16-разрядных чисел при сохранении
отдельных данных и операций в более точном представлении.
Дополнительные приёмы, включая масштабирование функции потерь,
учитывали ограниченный диапазон чисел~\cite{mixed2018}.
В этой работе исследовалось обучение. Для оценки энергии инференса необходимо
отдельно изучать выполнение уже обученной модели.

Смешанная точность предполагает назначение различных представлений
операндам, накопленным результатам и обновляемым параметрам.
Численная точность распределяется с учётом назначения операций
и условий устойчивости вычислительного процесса. При выборе форматов учитываются допустимые изменения результата, затраты
памяти и возможности вычислительной аппаратуры.

В методе HAWQ (Hessian AWare Quantization — квантование с учётом
матрицы вторых производных), предложенном Чжэнь Дуном и соавторами
в 2019 году, выбор точности связывался с различной
чувствительностью слоёв к квантованию~\cite{hawq2019}.
В этом подходе разрядность выбиралась для отдельных слоёв
с учётом их свойств. В методе HAQ
(Hardware-Aware Automated Quantization — автоматическое квантование
с учётом аппаратуры) К.~Вана и соавторов того же года использовалась
обратная связь от модели аппаратуры, включавшая оценки задержки
и энергии~\cite{haq2019}.
Выбор точности связывается с двумя характеристиками:
чувствительностью частей сети к изменению представления
и аппаратной стоимостью соответствующих численных режимов.

Разрядность можно выбирать при подготовке модели или во время её исполнения. В
первом случае автоматическая процедура может заранее задать постоянные форматы
для всех слоёв. Во втором случае система принимает решения по текущим данным и
должна выполнять необходимые управляющие действия.

\subsection{Округление и управление масштабом в работах 2010-х годов}

В работах 2010-х годов снижение разрядности исследовалось
совместно с правилами преобразования чисел и выбором масштаба.
Обучение ставило отдельную задачу: ограниченное представление
значений влияло на изменение параметров сети.
Поэтому численный формат рассматривался во взаимосвязи
с процедурой подготовки модели.

С.~Гупта и соавторы в 2015 году исследовали обучение
с 16-разрядной арифметикой с фиксированной точкой
и стохастическим округлением~\cite{gupta2015}.
Работа показала значение способа округления для обучения
с ограниченным представлением чисел. Исследование относилось
к подготовке параметров, тогда как развитие целочисленного
инференса было связано с исполнением уже обученной модели.
Эти задачи поддерживали разные направления исследований
низкоразрядных вычислений.

В Flexpoint, предложенном У.~Кёстером и соавторами
в 2017 году, значения одного тензора используют общий
показатель масштаба, изменяемый при исполнении~\cite{flexpoint2017}.
В центре этой работы находилась адаптация численного
представления к данным в процессе обучения.
Предложение проверялось на моделях нескольких нейросетевых
архитектур и рассматривалось как альтернатива использованию
обычного формата с плавающей точкой.

Работы Гупты и Кёстера показывают разные способы исследования
точности: через выбор правила округления и через изменение
масштаба представления. Обе постановки связывали численную
арифметику с особенностями нейросетевого обучения.
Таким образом, развитие форматов включало не только
сокращение количества битов, но и пересмотр способов
использования доступной точности.

\subsection{Форматы пониженной разрядности и блочное масштабирование}

Форматы пониженной разрядности разрабатывались не только для целых чисел.
В публикации \textit{FP8 Formats for Deep Learning} 2022 года
предложены два восьмиразрядных формата с плавающей точкой:
E4M3 и E5M2~\cite{fp82022}. Их различие отражало выбор между
диапазоном значений и точностью представления.
Работа развивала способы применения низкоразрядной арифметики
к обучению и использованию нейронных сетей, включая
языковые модели.

Появление таких предложений расширяло обсуждение численной
точности за пределы целочисленного инференса. Исследовались
разные способы представления чисел, отвечающие составу
операций и требованиям нейросетевых моделей.
Выбор формата становился частью совместного обсуждения
алгоритма обучения, качества результата и аппаратных затрат.

В работе Б.~Дарвиш Рухани и соавторов 2023 года исследованы форматы с
масштабированием небольших групп чисел (microscaling), в которых значения
пониженной разрядности используют общий коэффициент масштаба~\cite{mx2023}.
В таком представлении часть информации относится к целому блоку,
а не хранится заново в каждом его элементе. Размер блока,
формат масштаба и формат отдельного числа совместно определяют
поведение арифметики.

Выбор общей шкалы связывает статистику группы данных
с её машинным представлением.
Вместе с тем само определение формата ещё не задаёт устройство
детектора диапазона, порядок буферизации и способ выбора режима
при исполнении. Числовой формат задаёт представление данных, но способ его аппаратной
поддержки необходимо определять при проектировании устройства.
